\documentclass[11pt,a4paper]{article}
\usepackage[utf8]{inputenc}
\usepackage[T1]{fontenc}
\usepackage[polish]{babel}
\usepackage{lmodern}
\usepackage{geometry}
\geometry{margin=27mm}
\usepackage{amsmath}
\usepackage[hidelinks]{hyperref}

\title{Krótka edukacja o AI a niezależność poznawcza}
\author{Mateusz Pomianek}
\date{Projekt badania, wersja 0.9.0, 24 listopada 2026}

\begin{document}
\maketitle

\begin{abstract}
Skuteczne korzystanie z generatywnej sztucznej inteligencji nie musi prowadzić do poprawy zdolności samodzielnego rozumowania. Celem projektowanego badania jest oszacowanie wpływu krótkiej interwencji edukacyjnej, poświęconej ograniczeniom AI i weryfikacji informacji, na trafność niezależnego osądu. Proponuje się randomizowany eksperyment internetowy z aktywną grupą kontrolną, pomiarem bazowym, bezpośrednim posttestem oraz pomiarem odroczonym po siedmiu dniach. Głównym wynikiem będzie odsetek trafnych decyzji w nowych zadaniach wykonywanych bez pomocy AI. Dodatkowo ocenione zostaną kalibracja pewności, jakość uzasadnienia i wiedza o AI. Projekt pozwoli odróżnić przyrost wiedzy deklaratywnej od transferu kompetencji poznawczych, jednak nie umożliwi wnioskowania o globalnej zmianie inteligencji ani trwałej autonomii poza badanym kontekstem. Tekst przedstawia projekt badania, a nie wyniki empiryczne.
\end{abstract}

\noindent Słowa kluczowe: kompetencje AI, niezależność poznawcza, transfer uczenia się, automatyzacja osądu, kalibracja metapoznawcza.

\section{Wprowadzenie}
Systemy generatywnej AI umożliwiają szybkie uzyskanie odpowiedzi, ale płynność językowa nie stanowi dowodu poprawności wniosku. Użytkownik może uzyskać rozwiązanie zadania bez samodzielnej oceny przesłanek, co rodzi pytanie, czy doświadczenie korzystania z narzędzia przekłada się na późniejszą sprawność poznawczą bez jego wsparcia. W edukacji rozróżnienie to ma znaczenie zasadnicze: jakość wytworzonej odpowiedzi nie jest tożsama z kompetencją osoby, która tę odpowiedź przygotowała.

Niezależność poznawczą definiuje się tutaj wąsko jako zdolność do sformułowania trafnego osądu na podstawie dostępnych przesłanek bez bieżącej rekomendacji AI. Nie oznacza ona odrzucania narzędzi informatycznych, nieufności wobec wszystkich źródeł ani ogólnej cechy osobowości. Obejmuje rozpoznawanie brakujących informacji, rozróżnianie przesłanki od wniosku i dostosowanie pewności do jakości dowodów. Operacyjny pomiar tych umiejętności pozwala uniknąć utożsamienia deklarowanej samodzielności z rzeczywistym wykonaniem.

\section{Podstawy teoretyczne}
\subsection{Kompetencje AI i krytyczna ocena informacji}
Long i Magerko (2020) ujmują AI literacy jako zestaw kompetencji obejmujący rozumienie, krytyczną ocenę i skuteczne wykorzystywanie technologii. Chee i in. (2025) przedstawiają zróżnicowanie oczekiwanych kompetencji w zależności od etapu edukacyjnego oraz kontekstu zawodowego. W projekcie przyjmuje się, że znajomość ograniczeń modeli językowych jest potencjalnym warunkiem krytycznej oceny, lecz nie jest wystarczającym wskaźnikiem wykonania zadania. Badanie Le i in. (2026), dotyczące szkolenia przyszłych nauczycieli, wskazuje na znaczenie projektowania kompetencji wykraczających poza obsługę pojedynczego narzędzia. Jego wyniki dotyczą jednak m.in. poczucia własnej skuteczności i postaw; nie rozstrzygają przyjętego tutaj pytania o samodzielne wykonanie po usunięciu pomocy AI.

\subsection{Nadmierne poleganie na automatyzacji}
Badania nad automation bias opisują możliwość zastępowania własnej oceny rekomendacją systemu automatycznego (Skitka i in., 1999). Buçinca i in. (2021) wykazali, że wymuszenie aktywnego namysłu podczas decyzji wspomaganej przez AI może ograniczać nadmierne poleganie na błędnych wskazówkach. Nie wynika stąd jednak, że taki mechanizm utrzymuje się po zakończeniu interakcji. Interwencja edukacyjna może działać przez zmianę sposobu sprawdzania przesłanek, ale alternatywnym wyjaśnieniem krótkotrwałej poprawy jest większa ostrożność wynikająca z instrukcji eksperymentalnej. W preprincie Pupparta i Aru (2025) krótka interwencja nie ograniczyła istotnie przyjmowania błędnych rekomendacji ChatGPT, a zwiększyła odrzucanie trafnych porad. Badanie dotyczyło decyzji wspomaganych przez AI, nie samodzielnego posttestu. Wynik ten uzasadnia rozróżnienie trafności oceny od ogólnego sceptycyzmu: szkolenie nie powinno premiować odrzucania każdego twierdzenia, lecz poprawne rozpoznawanie, które twierdzenia wynikają z dowodów.

\subsection{Transfer uczenia się i mechanizm interwencji}
Barnett i Ceci (2002) podkreślają, że transfer należy określać względem podobieństwa treści, kontekstu, czasu i wymagań zadania. W proponowanym eksperymencie transfer oznacza zastosowanie reguł oceny dowodów do zadań niewystępujących w szkoleniu, a następnie ich ponowne zastosowanie po upływie tygodnia. Badanie Bastani i in. (2025) w nauczaniu matematyki uzasadnia oddzielanie wykonania wspomaganego od późniejszego wykonania bez AI. Nie stanowi ono jednak bezpośredniego dowodu skuteczności proponowanej interwencji u dorosłych.

Zakładany mechanizm polega na nabyciu procedury: sformułowanie problemu, identyfikacja przesłanek, poszukiwanie brakujących warunków i rozważenie alternatywnego wyjaśnienia przed wydaniem decyzji. Procedura ma być stosowana również wtedy, gdy rekomendacja AI nie jest dostępna. Rywalizujące wyjaśnienia obejmują efekt nowości, oczekiwania uczestnika, większy wysiłek i ogólny trening rozumowania. Randomizacja pozwoli oszacować efekt całego pakietu edukacyjnego, ale nie zidentyfikuje samodzielnie udziału każdego mechanizmu.

\section{Cel i hipotezy}
Pytanie główne brzmi: czy krótka edukacja o ograniczeniach AI i krytycznej weryfikacji poprawia trafność osądu w nowych zadaniach bez wsparcia AI w porównaniu z neutralną lekcją o porównywalnej długości?

\begin{enumerate}
\item H1 (pierwszorzędowa): uczestnicy interwencji uzyskają wyższy odsetek trafnych odpowiedzi w bezpośrednim postteście niż uczestnicy aktywnej kontroli, po uwzględnieniu wyniku bazowego.
\item H2 (drugorzędowa): przewaga interwencji będzie obserwowalna również w nowych zadaniach po siedmiu dniach.
\item H3 (drugorzędowa): interwencja zmniejszy średni błąd kwadratowy prognozy poprawności własnej odpowiedzi (wynik Briera). Kalibrację sensu stricto ocenimy dodatkowo na podstawie zależności między pewnością a obserwowaną trafnością.
\item H4 (drugorzędowa): uczestnicy interwencji uzyskają wyższy średni wynik w rubryce jakości uzasadnienia decyzji.
\end{enumerate}

Przyrost wiedzy o AI stanowi kontrolę manipulacji, nie dowód poprawy samodzielnego osądu. Związek między wiedzą, pewnością i wykonaniem będzie analizowany eksploracyjnie, bez automatycznego przypisywania zależnościom statusu mediacji przyczynowej.

\section{Materiał i metody}
\subsection{Schemat i uczestnicy}
Proponuje się dwuramienny eksperyment z alokacją 1:1 i pomiarem w trzech punktach czasowych. Populacją docelową są pełnoletnie osoby posługujące się językiem polskim, niezależnie od wcześniejszego doświadczenia z AI. Rekrutacja internetowa będzie dobrowolna. W przypadku rekrutacji studentów udział nie może mieć wpływu na ocenę z przedmiotu ani pozostawać w zależności od decyzji prowadzącego.

Przed randomizacją uczestnicy wykonają test bazowy i odpowiedzą na krótkie pytania o korzystanie z AI. Przydział zostanie wygenerowany komputerowo i ujawniony dopiero po ukończeniu pretestu. Proponowany jest losowy przydział warstwowy względem dwóch poziomów doświadczenia z AI, z progami zdefiniowanymi przed rekrutacją. Uczestnicy nie mogą być w pełni zaślepieni na treść lekcji; nie zostaną jednak poinformowani o oczekiwanym kierunku różnic. Osoby oceniające uzasadnienia nie będą znały grupy ani punktu pomiarowego.

\subsection{Interwencja i kontrola}
Interwencja będzie trwała około 15 minut i obejmie trzy części: wyjaśnienie różnicy między przekonującym tekstem a uzasadnionym twierdzeniem, ćwiczenia rozpoznawania brakujących przesłanek oraz zastosowanie czteroetapowej procedury weryfikacji. Uczestnik przed ujawnieniem komentarza dydaktycznego podejmie własną decyzję. Przykłady błędnych odpowiedzi AI zostaną przygotowane wcześniej, zweryfikowane merytorycznie i będą identyczne dla wszystkich osób w tym ramieniu.

Kontrola aktywna obejmie lekcję o historii i podstawowych zastosowaniach narzędzi cyfrowych, bez instruktażu krytycznej oceny dowodów. Robocze lekcje obejmują po dziesięć ekranów i cztery pytania z informacją zwrotną po odpowiedzi. Zakłada się taki sam budżet czasowy i interfejs; porównywalność rzeczywistego wysiłku, czasu i zainteresowania wymaga pilotażu. Dwa ramiona pozwalają ocenić całkowity efekt szkolenia; wyodrębnienie komponentu specyficznego dla AI względem ogólnego treningu krytycznego myślenia wymagałoby dodatkowego ramienia i osobnego planu mocy.

\subsection{Bodźce i procedura pomiarowa}
Przygotowano roboczy bank pilotażowy trzech zestawów po 24 zadania, obejmujące ocenę wnioskowania, brakującego kontekstu i niespójności dowodów. W każdym zadaniu uczestnik wskaże jedną z trzech odpowiedzi: wniosek uzasadniony, wniosek sprzeczny z przesłankami lub brak wystarczających informacji. Po decyzji poda subiektywne prawdopodobieństwo jej poprawności w zakresie od 0 do 100 procent. W wybranych sześciu zadaniach każdego zestawu poda również krótkie uzasadnienie. Wszystkie informacje konieczne do odpowiedzi muszą znajdować się w treści zadania. W każdej formie przewiduje się osiem pozycji z każdą z trzech poprawnych kategorii odpowiedzi. Treść zostanie rozłożona na trzy obszary: codzienne decyzje, interpretacja prostych danych oraz ocena twierdzeń o zależnościach. Zbilansowanie kategorii ograniczy możliwość uzyskania wysokiego wyniku przez systematyczne wybieranie odpowiedzi o braku informacji. Kolejność odpowiedzi będzie losowana, a znaczenie każdej kategorii zostanie wyjaśnione w identycznej instrukcji przed pretestem.

Przykład pozycji treningowej: w dwóch grupach odnotowano różne średnie wyniki testu, ale nie podano sposobu przydziału uczestników ani wyników bazowych. Wniosek, że różnica została spowodowana wyłącznie przez zastosowaną metodę nauczania, nie jest uzasadniony; prawidłowa kategoria to brak wystarczających informacji. Dla odróżnienia pozycji sprzecznych od nierozstrzygalnych klucz musi wskazywać konkretną przesłankę obalającą twierdzenie albo brak przesłanki koniecznej do rozstrzygnięcia. Przykład ten nie wejdzie do zestawów testowych.

Zestawy będą równoważone między punktami pomiarowymi i grupami według sześciu możliwych kolejności. Ich treści nie będą powtarzać przykładów szkoleniowych. Niezależni eksperci ocenią zgodność klucza, jednoznaczność i trudność, a pilotaż z osobami spoza próby głównej umożliwi wykrycie efektów sufitowych, podłogowych i nierównoważnych pozycji. Bank roboczy zawiera klucze i jest publicznie dostępny; nie będzie używany bez zmian jako test główny. Ostateczne formy zostaną przygotowane z kontrolą dostępu i poddane odrębnej ocenie oraz pilotażowi. Ostateczna liczba pozycji i czas odpowiedzi zostaną zamrożone po pilotażu, przed rozpoczęciem rekrutacji głównej. Proponuje się ocenę treści przez co najmniej dwóch niezależnych ekspertów oraz pilotaż poznawczy z 8--12 osobami, a następnie pilotaż ilościowy z około 60 osobami. Progi robocze to zgodność ekspertów w zakresie klucza dla wszystkich zachowanych pozycji, trudność w przedziale 0,25--0,90 i dodatnia skorygowana korelacja pozycji z wynikiem całkowitym. Przekroczenie progu jest przesłanką do oceny treści i ponownego pilotażu, nie automatycznego usunięcia pozycji. Rzetelność i różnice trudności między formami należy raportować z przedziałami niepewności. Pilotaż nie będzie wykorzystany jako test H1.

Pomiar bazowy i bezpośredni posttest odbędą się w jednej sesji internetowej. Pomiar odroczony nastąpi po siedmiu dniach, z dopuszczalnym oknem od sześciu do ośmiu dni. W żadnym z trzech testów nie będzie dostępna wbudowana pomoc AI ani informacja o poprawności odpowiedzi. Uczestnik zostanie poproszony o niewykorzystywanie zewnętrznych narzędzi. Oświadczenie o przestrzeganiu instrukcji jest wskaźnikiem pomocniczym, nie gwarancją braku pomocy.

\subsection{Zmienne i operacjonalizacja}
Pierwszorzędowy wskaźnik: proporcja poprawnych odpowiedzi na pozycje bezpośredniego posttestu. Dla osoby \(i\) w pomiarze \(t\):
\[
Y_{it}=\frac{1}{|O_{it}|}\sum_{j\in O_{it}} y_{itj},
\qquad y_{itj}\in\{0,1\},
\]
gdzie \(O_{it}\) oznacza zbiór pozycji z udzieloną odpowiedzią, a \(y_{itj}=1\) odpowiedź poprawną. Wynik oblicza się przy co najmniej 20 z 24 odpowiedzi; w przeciwnym przypadku jest brakujący. Zestaw nieukończony nie jest automatycznie klasyfikowany jako zbiór błędnych odpowiedzi. Zasady postępowania z brakami opisano w planie analizy.

Trafność probabilistycznej oceny własnej odpowiedzi zostanie oceniona za pomocą binarnego wyniku Briera (Brier, 1950):
\[
B_{it}=\frac{1}{|C_{it}|}\sum_{j\in C_{it}}(p_{itj}-y_{itj})^2,
\]
gdzie \(C_{it}\) oznacza zbiór pozycji z kompletną parą odpowiedź--pewność, a \(p_{itj}\) to deklarowane prawdopodobieństwo poprawności wybranej odpowiedzi w skali od 0 do 1. Wymaga się co najmniej 20 kompletnych par. Niższy wynik oznacza mniejszy błąd prognozy. Ponieważ miara uwzględnia zarówno trafność, jak i pewność, dodatkowo przedstawione zostaną średnia pewność, trafność i wykres kalibracji. Nie będzie ona traktowana jako czysta miara metapoznania. Dodatkowo obliczymy średnią różnicę między deklarowaną pewnością a trafnością, aby rozróżnić nadmierną i niedostateczną pewność. Osobno zostaną przedstawione macierze pomyłek dla trzech kategorii odpowiedzi. Wzrost częstości odpowiedzi o braku informacji kosztem poprawnego rozpoznawania uzasadnionych wniosków będzie traktowany jako możliwy efekt uboczny interwencji, a nie poprawa niezależności poznawczej.

Jakość argumentacji zostanie oceniona według trzystopniowej rubryki: 0 oznacza brak istotnej przesłanki, 1 wskazanie przesłanki bez poprawnego powiązania z wnioskiem, a 2 poprawne powiązanie przesłanki z oceną wniosku. Osoby oceniające zastosują ustalony wcześniej klucz i rozstrzygną niezgodności po niezależnym kodowaniu. Proponuje się raportowanie ważonej kappy dla ocen przed uzgodnieniem. Brak uzasadnienia nie jest automatycznie równoważny ocenie 0; średni wynik zostanie obliczony przy co najmniej pięciu z sześciu ocenionych uzasadnień.

\subsection{Liczebność i wykonalność}
Wielkość próby zostanie określona na podstawie najmniejszej różnicy mającej znaczenie edukacyjne, a nie efektu zaobserwowanego w małym pilotażu. Na etapie planowania przyjmuje się roboczo efekt standaryzowany \(d=0{,}35\), moc 80 procent i dwustronny poziom istotności \(0{,}05\). Przy założeniu niezależnych grup orientacyjne przybliżenie normalne daje około 129 osób na grupę. Przy założonej utracie 15 procent uczestników orientacyjny cel rekrutacji wynosi 304 osoby. To wyłącznie scenariusz planistyczny, który nie wykorzystuje jeszcze korelacji z pretestem. Ostateczne obliczenie uwzględni rzetelność zestawów, oczekiwaną korelację pomiarów, brak danych oraz model główny. Liczebność nie będzie ustalana na podstawie obserwowanych wartości \(p\).

Przy założeniu wspólnego nachylenia regresji i jednakowej wariancji reszt przybliżenie ANCOVA skaluje wymaganą liczebność przez czynnik \(1-\rho^2\), gdzie \(\rho\) oznacza korelację pretestu i posttestu. Dla \(d=0{,}35\) oraz 15-procentowej utraty uczestników scenariusze \(\rho=0\), \(0{,}3\), \(0{,}5\) i \(0{,}7\) odpowiadają rekrutacji odpowiednio 304, 276, 230 i 156 osób. Są to wartości modelowe, a nie oszacowania z pilotażu. Nie należy przyjmować optymistycznej korelacji bez uzasadnienia. W analizie wrażliwości rozważymy ponadto efekty mniejsze od zakładanego i różne odsetki braków danych.

W pilotażu dla każdej formy obliczymy odsetki poprawnych odpowiedzi z przedziałami Wilsona, skorygowane korelacje pozycji z wynikiem pozostałych pozycji oraz współczynnik KR-20 z bootstrapowym przedziałem niepewności. Trudność będzie obliczana wśród dostępnych odpowiedzi, a rzetelność i dyskryminacja na kompletnych formularzach, z jawnym raportowaniem liczebności i braków. Porównania form u tych samych osób uwzględnią kolejność ich prezentacji opisowo. Brak istotnej różnicy ani wysoka korelacja nie dowodzą równoważności form; formalna ocena wymaga ustalenia marginesu równoważności i odpowiedniej liczebności.

\section{Plan analizy}
\subsection{Estymand i analiza pierwszorzędowa}
Populacyjny estymand stanowi różnica średniej proporcji poprawnych odpowiedzi bezpośrednio po interwencji między warunkami losowego przydziału. Interesuje nas efekt przydziału do szkolenia, niezależnie od stopnia zaangażowania w materiał. Główna analiza wykorzysta regresję liniową z korektą wyniku bazowego (Vickers i Altman, 2001):
\[
Y_{i1}=\beta_0+\tau A_i+\beta_1Y_{i0}
+\boldsymbol{\gamma}^{\mathsf T}\mathbf{X}_i+\varepsilon_i,
\]
gdzie \(A_i\) oznacza przydział do interwencji, a \(\mathbf X_i\) obejmuje warstwę randomizacji i kolejność zestawów. Parametr \(100\tau\) będzie raportowany jako różnica w punktach procentowych wraz z 95-procentowym przedziałem ufności opartym na odpornych błędach standardowych HC3. W implementacji referencyjnej stosuje się asymptotyczne przedziały normalne; ich własności i ewentualną korektę dla małych prób należy zweryfikować przed zamrożeniem protokołu. H1 zostanie oceniona testem dwustronnym przy \(\alpha=0{,}05\), z oczekiwanym dodatnim kierunkiem efektu.

W analizie wrażliwości proponuje się model logistyczny na poziomie obserwowanych pozycji z losowym wyrazem wolnym uczestnika i pozycji, z korektą pretestu, warstwy i kolejności form. Referencyjna implementacja wykorzystuje bayesowskie przybliżenie wariacyjne; wymaga porównania z niezależnym dopasowaniem i oceny wrażliwości na rozkłady a priori. Efekt na skali prawdopodobieństwa zostanie przedstawiony przez standaryzację predykcji względem uczestników z dostępnym pretestem i całkowanie efektów losowych. Przedział wiarygodności z przybliżonego rozkładu a posteriori nie jest tożsamy z przedziałem ufności ANCOVA. Analiza pozycji nie zastępuje analizy zgodnej z przydziałem przy brakach danych. Rozbieżność między modelami będzie raportowana, a nie rozstrzygana wyborem korzystniejszego wyniku.

\subsection{Analizy drugorzędowe i brakujące dane}
H2 zostanie oceniona analogicznym modelem dla pomiaru odroczonego. H3 i H4 wykorzystają odpowiednio wynik Briera i średnią punktację uzasadnienia w postteście, z korektą odpowiedniej miary bazowej. W obrębie trzech testów drugorzędowych zastosowana zostanie korekta Holma. Interakcja z doświadczeniem w korzystaniu z AI oraz analiza wiedzy o AI pozostaną eksploracyjne. Nie należy interpretować nieistotnego statystycznie efektu jako dowodu braku różnicy.

Analiza obejmie wszystkich losowo przydzielonych uczestników zgodnie z przydziałem. Nieukończone pomiary, przyczyny rezygnacji i odsetki braków zostaną opisane według grup. Przy brakach wyniku zostanie zastosowana wielokrotna imputacja na poziomie wyników zestawów, z uwzględnieniem grupy, pretestu, kolejności zestawów i dostępnych pomiarów, przy założeniu braków losowych warunkowo względem obserwowanych danych. Proponuje się 50 imputacji z dopasowaniem średniej predykcyjnej (predictive mean matching), łączeniem estymat według reguł Rubina i kontrolą stabilności estymat. Wynik zostanie obliczony przy co najmniej 20 odpowiedziach z 24; przy mniejszej liczbie cały wynik pomiaru będzie brakujący. W analizie wrażliwości sprawdzimy wyłącznie kompletne zestawy oraz skrajne kodowania brakujących pozycji. Wrażliwość na naruszenie tego założenia zostanie oceniona przez analizę kompletnych przypadków oraz przesunięcie imputowanych wyników w kierunku gorszego wykonania w każdym ramieniu. Przed rekrutacją główną należy zamrozić te reguły punktowania, model imputacji i wartości przesunięć w analizie wrażliwości. Wstępnie proponuje się zmniejszenie imputowanej trafności bezpośredniego posttestu o 5 lub 10 punktów procentowych, osobno dla każdego ramienia, z zachowaniem ograniczenia skali i niezmienionych wartości obserwowanych. Przygotowano referencyjne obliczenia ANCOVA, korekty Holma i reguł Rubina dla dostarczonych imputacji; przygotowano również referencyjny generator PMM z równaniami łańcuchowymi, pięcioma najbliższymi dawcami z tego samego ramienia, 20 iteracjami na niezależnie inicjowany łańcuch i rejestracją przebiegu imputowanych średnich oraz odchyleń standardowych. Parametry te są propozycją, a zbieżność, stabilność względem nasiona i liczby iteracji, własności w warunkach braków oraz model pozycji wymagają dalszej weryfikacji. Analiza osób deklarujących przestrzeganie instrukcji będzie dodatkowa, gdyż wykluczenie po randomizacji może wprowadzać obciążenie.

\subsection{Kontrola jakości i przejrzystość}
Zamrożenie estymandu, reguł punktowania i kryteriów jakości nastąpi przed akwizycją główną. Należy ustalić procedurę postępowania z duplikatami oraz reguły oceny nadmiernie szybkiego odpowiadania. Sama szybkość nie może być podstawą arbitralnego usuwania uczestnika. Należy ocenić rzetelność zestawów i równoważność form w pilotażu oraz podać niepewność tych oszacowań. Skrypty analizy powinny zostać przetestowane na danych syntetycznych przed otwarciem zbioru głównego. Dane syntetyczne służą wyłącznie kontroli obliczeń. Opracowano generator ograniczonych wyników z mechanizmami MCAR, MAR i MNAR, który pozwala porównać analizę pełnych danych, kompletnych przypadków, PMM i przesunięcia imputowanych wartości. Plan obejmuje obciążenie, RMSE, pokrycie przedziałów, odsetek odrzuceń i nieudane dopasowania. Symulacja maskuje na razie tylko bezpośredni wynik trafności; nie reprezentuje pełnej rezygnacji uczestnika. Własności statystyczne wymagają wielokrotnych replikacji, a nie pojedynczego testu technicznego. Zaplanowano macierz 24 scenariuszy obejmującą trzy mechanizmy braków, dwie liczebności, dwa efekty populacyjne i dwa poziomy prawdopodobieństwa rezygnacji. Raporty techniczne i scenariusze rozpoznawcze będą oddzielone od pełnych symulacji. Przed zamrożeniem projektu należy uzasadnić minimalny efekt edukacyjny, zweryfikować mechanizmy braków całej sesji i ustalić ostateczną liczebność. Projekt prerejestracji przygotowano, lecz nie zarejestrowano.

\section{Etyka i ochrona danych}
Rekrutacja rozpocznie się dopiero po uzyskaniu właściwej oceny etycznej i zatwierdzeniu informacji dla uczestnika. Uczestnicy otrzymają opis dobrowolności, czasu udziału, możliwości rezygnacji i zasad przetwarzania danych. Kontakt potrzebny do pomiaru odroczonego powinien być przechowywany odrębnie od odpowiedzi, z ograniczonym dostępem i okresem retencji. Uczestnikom nie należy zlecać wprowadzania danych osobowych do zewnętrznych usług AI. Materiały dydaktyczne nie będą zawierać niebezpiecznych zaleceń ani rzeczywistych danych osób trzecich. Udostępnienie materiałów i danych będzie przedmiotem odrębnej decyzji, zależnej od zgód i ryzyka identyfikacji.

\section{Wyniki}
Nie rozpoczęto akwizycji danych na potrzeby opisanego eksperymentu. Nie oszacowano efektu, nie przeprowadzono testów hipotez i nie potwierdzono wykonalności procedury. Po zakończeniu badania sekcja obejmie przepływ uczestników, charakterystykę próby, statystyki opisowe, oszacowania efektów oraz wyniki analiz wrażliwości.

\section{Dyskusja i ograniczenia}
Potwierdzenie H1 wskazywałoby, że szkolenie może zwiększyć trafność decyzji bez bieżącej rekomendacji AI w obrębie zdefiniowanych zadań. Potwierdzenie H2 dostarczyłoby dodatkowej informacji o utrzymaniu efektu przez tydzień. Poprawa jedynie wiedzy deklaratywnej lub wzrost pewności bez wzrostu trafności nie uzasadniałaby wniosku o poprawie niezależności poznawczej. W przypadku braku oczekiwanej różnicy należy rozważyć zarówno nieskuteczność interwencji, jak i nieprecyzyjny pomiar lub niedostateczny transfer.

Internetowy charakter badania ogranicza kontrolę nad korzystaniem z pomocy zewnętrznej. Równoważenie form nie eliminuje całkowicie różnic trudności, a rekrutacja ochotników może ograniczać reprezentatywność próby. Uczestnicy mogą odgadnąć cel szkolenia, co może zwiększyć staranność niezależnie od trwałej zmiany kompetencji. Zadania zamknięte redukują koszt i obciążenie oceną, lecz słabiej reprezentują złożone decyzje w realnym środowisku. Siedmiodniowy pomiar nie pozwala wnioskować o wielomiesięcznej trwałości efektu. Odmienna treść zadań przy podobnej strukturze rozumowania pozwoli badać transfer do nowych pozycji, ale nie wystarcza do wykazania dalekiego transferu. Jawność roboczego banku stanowi dodatkowe ryzyko znajomości odpowiedzi; zamrożone zestawy główne muszą być odrębne i niedostępne uczestnikom przed pomiarem.

\section{Wnioski}
Projekt proponuje behawioralny test edukacji AI ukierunkowanej na samodzielne rozumowanie. Wkładem ma być oddzielenie kompetencji ocenianej bez pomocy AI od deklarowanej wiedzy i jakości pracy wspomaganej. Przed rozpoczęciem badań konieczne są pilotaż narzędzia, uzasadnienie docelowej liczebności, zamrożenie protokołu i spełnienie wymagań etycznych. Na obecnym etapie nie można ocenić skuteczności interwencji.

\section*{Bibliografia}
\begingroup
\setlength{\parindent}{0pt}
\setlength{\parskip}{0.5em}
\hangindent=1.5em Barnett, S. M., \& Ceci, S. J. (2002). When and where do we apply what we learn? A taxonomy for far transfer. \textit{Psychological Bulletin, 128}(4), 612--637. \url{https://doi.org/10.1037/0033-2909.128.4.612}

\hangindent=1.5em Bastani, H., Bastani, O., Sungu, A., Ge, H., Kabakc{\i}, Ö., \& Mariman, R. (2025). Generative AI without guardrails can harm learning: Evidence from high school mathematics. \textit{Proceedings of the National Academy of Sciences, 122}(26), e2422633122. \url{https://doi.org/10.1073/pnas.2422633122}

\hangindent=1.5em Brier, G. W. (1950). Verification of forecasts expressed in terms of probability. \textit{Monthly Weather Review, 78}(1), 1--3. \url{https://doi.org/10.1175/1520-0493%281950%29078%3C0001%3AVOFEIT%3E2.0.CO%3B2}

\hangindent=1.5em Buçinca, Z., Malaya, M. B., \& Gajos, K. Z. (2021). To trust or to think: Cognitive forcing functions can reduce overreliance on AI in AI-assisted decision-making. \textit{Proceedings of the ACM on Human-Computer Interaction, 5}(CSCW1), Article 188. \url{https://doi.org/10.1145/3449287}

\hangindent=1.5em Chee, H., Ahn, S., \& Lee, J. (2025). A competency framework for AI literacy: Variations by different learner groups and an implied learning pathway. \textit{British Journal of Educational Technology, 56}(5), 2146--2182. \url{https://doi.org/10.1111/bjet.13556}

\hangindent=1.5em Le, T. H., Huynh, L., Dang, B., Pham, H.-H., Nguyen, N. T. V., \& Nguyen, A. (2026). Development and evaluation of artificial intelligence literacy training for teacher education students. \textit{British Journal of Educational Technology}. Advance online publication. \url{https://doi.org/10.1111/bjet.70047}

\hangindent=1.5em Long, D., \& Magerko, B. (2020). What is AI literacy? Competencies and design considerations. In \textit{Proceedings of the 2020 CHI Conference on Human Factors in Computing Systems} (pp. 1--16). Association for Computing Machinery. \url{https://doi.org/10.1145/3313831.3376727}

\hangindent=1.5em Puppart, B., \& Aru, J. (2025). \textit{Short-term AI literacy intervention does not reduce over-reliance on incorrect ChatGPT recommendations} [Preprint]. arXiv. \url{https://arxiv.org/abs/2503.10556}

\hangindent=1.5em Skitka, L. J., Mosier, K. L., \& Burdick, M. (1999). Does automation bias decision-making? \textit{International Journal of Human-Computer Studies, 51}(5), 991--1006. \url{https://doi.org/10.1006/ijhc.1999.0252}
\hangindent=1.5em Vickers, A. J., \& Altman, D. G. (2001). Analysing controlled trials with baseline and follow up measurements. \textit{BMJ, 323}(7321), 1123--1124. \url{https://doi.org/10.1136/bmj.323.7321.1123}
\endgroup
\end{document}
